% Self-contained section fragment; requires amsmath, amssymb, amsthm.
% The parent document supplies theorem, lemma, proposition, corollary, remark.
% Labels and bibliography keys in this contribution have the prefix sd:.
\section{Higher spectral derivatives through the faces of a simplex}
\label{sd:sec:main}

The incoming beta-polarization report asks for explicit higher-order
subtractions on a Dirichlet simplex and for their extension to odd
polynomial weights \cite[Further research, items 3--4]{sd:HigherDefects}.
The incoming quartic report already gives a finite-subtraction compiler
at every spectral order \cite[Section 5]{RepoQuartic}.  We give a different
answer: a finite representation by ordinary simplex integrals of
\emph{lower} one-mass derivatives.  Its useful structural feature is that
an interaction involving a face with $k$ vertices cannot appear before
the $k$th spectral derivative, after an explicitly computable zero-jet
polynomial has been separated.

The scalar Dirichlet parameter identity is classical.  It is already
used for mixed twists in the canonical manuscript
\cite[\texttt{07-mixed-twist-identities.tex}]{RepoMain}; logarithmic
Dirichlet averages have a substantial classical theory
\cite{sd:Carlson1987,sd:Carlson1991}.  The contribution here is the explicit
face subtraction, its absolute-convergence proof, the resulting
all-order formula, and the rational zero-jet term.  This is a claim of
new content relative to the inspected repository, not an assertion of
exhaustive literature priority.

\subsection{Spectral normalization and the zero-jet polynomial}
\label{sd:subsec:setup}

Fix $a>0$, $r\ge2$, real masses $B_1,\ldots,B_r\ge0$, and a polynomial
$W(x)=\sum_{d=0}^{D}w_dx^d$.  Write $x_n=n+a$ and initially define
\begin{align}
 E_W(s;B)&=\sum_{n\ge0}W(x_n)(x_n^2+B)^{-s},
       \label{sd:eq:one}\\
 Z_{W,\boldsymbol B}(s)&=\sum_{n\ge0}W(x_n)
                    \prod_{i=1}^r(x_n^2+B_i)^{-s},\label{sd:eq:many}\\
 \mathcal D_W(s)&=\frac1r\sum_{i=1}^rE_W(rs;B_i)
                          -Z_{W,\boldsymbol B}(s),\qquad
 \Delta_m(W)=\mathcal D_W^{(m)}(0).
       \label{sd:eq:defect}
\end{align}
The sums start in their half-planes of absolute convergence and are then
continued meromorphically.  All logarithms of the positive quadratic
factors are real.  Let
\begin{equation}\label{sd:eq:jets}
 F_j(B)=\left.\partial_s^j E_W(s;B)\right|_{s=0}.
\end{equation}

\begin{lemma}[The polynomial at spectral order zero]
\label{sd:lem:zero}
The spectral functions in \eqref{sd:eq:one}--\eqref{sd:eq:many} are
holomorphic at zero.  The one-mass zero jet is
\begin{equation}\label{sd:eq:zero}
 F_0(B)=\sum_{d=0}^{D}w_d\zeta(-d,a)
           +\sum_{h\ge1}p_h B^h,\qquad
 p_h=\frac{(-1)^h}{2h}\,w_{2h-1},
\end{equation}
where missing polynomial coefficients are zero.  In particular $F_0$
is mass-independent for even $W$.
\end{lemma}
\begin{proof}
For sufficiently large $x$, expand the quadratic factor with a finite
remainder:
\[
 (x^2+B)^{-s}=x^{-2s}\sum_{h=0}^{K}
                 \frac{(-1)^h(s)_h}{h!}B^h x^{-2h}
                  +\text{remainder}.
\]
After multiplication by $x^d$, the removed terms sum to
$(-1)^h(s)_hB^h\zeta(2s+2h-d,a)/h!$.  Choose $K$ large enough
that the remainder and its required derivatives converge normally near
zero.  A finite head is retained exactly, so no restriction such as
$B<a^2$ is needed.  For $h\ge1$, the coefficient vanishes at zero
unless $2h-d=1$.  In that case
\[
 \frac{(s)_h}{h!}=\frac{s}{h}+O(s^2),\qquad
 \zeta(1+2s,a)=\frac1{2s}+O(1),
\]
which gives the term $(-1)^h B^h/(2h)$ in \eqref{sd:eq:zero}.
The same subtraction for the product of quadratic factors has positive
Taylor-degree coefficients vanishing at $s=0$; they cancel its possible
simple Hurwitz poles.  This also proves holomorphy of the product
spectrum.  The argument is locally uniform in the masses and in $a>0$.
\end{proof}

\subsection{An explicit subtraction on every face}
\label{sd:subsec:faces}

For a nonempty set $S\subseteq\{1,\ldots,r\}$, let
\[
 \Sigma_S=\{(t_i)_{i\in S}:t_i\ge0,\ \sum_{i\in S}t_i=1\}.
\]
The measure is ordinary $(|S|-1)$-dimensional coordinate measure:
$dt_1\cdots dt_{k-1}$ after the final coordinate has been eliminated.
Thus the simplex has volume $1/(k-1)!$, where $k=|S|$.
For a nonempty $T\subseteq S$ put
\[
 t_T=\sum_{i\in T}t_i,\qquad
 B_T=\frac{\sum_{i\in T}t_i B_i}{t_T},\qquad
 G_T(F;t)=t_TF(B_T).
\]
Set $G_T=0$ when $t_T=0$.  Define the face subtraction
\begin{equation}\label{sd:eq:H}
 \boxed{\displaystyle
 \mathcal H_S(F;t)=
  \sum_{\varnothing\ne T\subseteq S}
       (-1)^{|S|-|T|}\,t_TF(B_T).}
\end{equation}
For a pair this is the difference from the affine chord:
\begin{equation}\label{sd:eq:Hpair}
 \mathcal H_{ij}(F;u)
   =F(uB_i+(1-u)B_j)-uF(B_i)-(1-u)F(B_j).
\end{equation}
For three vertices it is
\begin{align}\label{sd:eq:Htriple}
 \mathcal H_{ijk}(F;t)={}&F(t_iB_i+t_jB_j+t_kB_k)
                  +\sum_{\ell\in\{i,j,k\}}t_\ell F(B_\ell)\notag\\
 &-\sum_{\{\ell,m\}\subseteq\{i,j,k\}}
   (t_\ell+t_m)F\!\left(
            \frac{t_\ell B_\ell+t_mB_m}{t_\ell+t_m}\right).
\end{align}
These formulas remain regular when masses coincide.  They use neither
division by a mass difference nor an unspecified finite part.

\begin{lemma}[Boundary cancellation and ordinary integrability]
\label{sd:lem:boundary}
If $F\in C^1([\min B_i,\max B_i])$, then $\mathcal H_S$ vanishes
on every proper coordinate face and
\begin{equation}\label{sd:eq:bound}
 |\mathcal H_S(F;t)|\le C\min_{i\in S}t_i.
\end{equation}
For $k=|S|\ge2$, the integral
\begin{equation}\label{sd:eq:integralhol}
 \int_{\Sigma_S}\mathcal H_S(F;t)
                            \prod_{i\in S}t_i^{s-1}\,dt
\end{equation}
is absolutely convergent for $\operatorname{Re}s>-1/(k-1)$.
It is holomorphic there, with locally normally convergent fixed
logarithmic derivatives.  The same conclusions hold for a holomorphic
family of $F$ whose $C^1$ norms are locally bounded.
\end{lemma}
\begin{proof}
The perspective $G_T=t_TF(B_T)$ is Lipschitz on the nonnegative
coordinate cone.  Away from $t_T=0$,
\[
 \partial_{t_i}G_T=F(B_T)+F'(B_T)(B_i-B_T)\quad(i\in T),
\]
so its coordinate derivatives have a uniform bound.  Continuity at the
origin extends the Lipschitz bound there.  If $t_i=0$, the terms indexed
by $T$ and $T\cup\{i\}$ cancel in \eqref{sd:eq:H}.  Changing that
coordinate from zero to $t_i$ therefore gives a bound $Ct_i$ for
every $i$, proving \eqref{sd:eq:bound}.

Cover the simplex by the $k$ regions where a specified coordinate is
at least $1/k$.  On such a region,
\[
 \min_i t_i\le
        \left(\prod_{i\ne i_{\max}}t_i\right)^{1/(k-1)}.
\]
The remaining coordinate factors are dominated by a product of
integrable powers whenever $\operatorname{Re}s+1/(k-1)>0$.
Any fixed power of $\sum_i|\log t_i|$ remains integrable on smaller
closed half-planes.  Dominated holomorphic parameter integration proves
all assertions.  In particular, every integral below with the weight
$\prod_i t_i^{-1}$ is an ordinary absolutely convergent integral.
\end{proof}

\begin{theorem}[Exact decomposition by simplex faces]
\label{sd:thm:master}
Put
\begin{equation}\label{sd:eq:Q}
 Q_k(s)=\frac{\Gamma(1+ks)}{\Gamma(1+s)^k}.
\end{equation}
The following identity holds meromorphically for
$\operatorname{Re}s>-1/(r-1)$, with ordinary face integrals away
from spectral poles, and in particular as a holomorphic germ at zero:
\begin{equation}\label{sd:eq:master}
 \boxed{\displaystyle
 \mathcal D_W(s)=-\frac1r\sum_{k=2}^{r}s^{k-1}Q_k(s)
   \sum_{|S|=k}\int_{\Sigma_S}
     \mathcal H_S(E_W(rs;\cdot);t)
                         \prod_{i\in S}t_i^{s-1}\,dt.}
\end{equation}
\end{theorem}
\begin{proof}
The product of Gamma integrals, with variables changed to a radial
coordinate and simplex proportions, gives
\[
 \prod_{i=1}^r y_i^{-s}
 =\frac{\Gamma(rs)}{\Gamma(s)^r}
   \int_{\Sigma_{\{1,\ldots,r\}}}
        \left(\sum_i t_i y_i\right)^{-rs}
                            \prod_i t_i^{s-1}\,dt
\]
for $y_i>0$ and $\operatorname{Re}s>0$.  Apply this to
$y_i=x_n^2+B_i$ and sum where the spectra converge absolutely.
Finite subset inversion gives
\[
 F\!\left(\sum_i t_iB_i\right)
       =\sum_{\varnothing\ne S\subseteq\{1,\ldots,r\}}
                    \mathcal H_S(F;t).
\]
On the full simplex, each $\mathcal H_S$ is homogeneous of degree one
in the coordinates belonging to $S$.  Under the normalized symmetric
Dirichlet weight, $q=t_S$ has mean $|S|/r$ and is independent of the
normalized proportions $t_i/q$, $i\in S$.  Equivalently, this follows
by integrating first in $q$ and the complementary proportions using
the ordinary beta integral.  Therefore, for $k=|S|$,
\[
 \mathbb E_{r,s}\mathcal H_S(F;t)
 =\frac{k}{r}\frac{\Gamma(ks)}{\Gamma(s)^k}
       \int_{\Sigma_S}\mathcal H_S(F;t)
                                 \prod_{i\in S}t_i^{s-1}\,dt.
\]
The singleton terms are $r^{-1}\sum_iF(B_i)$.  Subtracting them and
using
$\frac{k}{r}\Gamma(ks)/\Gamma(s)^k=s^{k-1}Q_k(s)/r$
proves \eqref{sd:eq:master} in the initial half-plane.
The spectral finite subtractions in Lemma~\ref{sd:lem:zero}, together
with Lemma~\ref{sd:lem:boundary}, give meromorphic continuation of the
integrals and the asserted holomorphic germs.  The equality follows
there by continuation from the common convergence domain.
\end{proof}

\subsection{The entire zero-jet contribution is a finite rational function}
\label{sd:subsec:local}

Separate $E_W(rs;B)=F_0(B)+(E_W(rs;B)-F_0(B))$ in
\eqref{sd:eq:master}.  Only the mass polynomial in
\eqref{sd:eq:zero} contributes to the first part.

\begin{theorem}[Rational zero-jet contribution]
\label{sd:thm:local}
The part of $\mathcal D_W(s)$ arising from $F_0$ is
\begin{equation}\label{sd:eq:local}
 \boxed{\displaystyle
 \mathcal L_W(s)=\sum_{h\ge1}p_h\left\{
 \frac1r\sum_{i=1}^r B_i^h-
 \frac{h!}{(rs)_h}
       \sum_{h_1+\cdots+h_r=h}
            \prod_{i=1}^r\frac{(s)_{h_i}B_i^{h_i}}{h_i!}
                                      \right\}.}
\end{equation}
All apparent singularities at $s=0$ are removable, and
$\mathcal L_W(0)=0$.  The term $h=1$ vanishes identically.  For
$m\ge1$, set
\begin{equation}\label{sd:eq:Ccoeff}
 C_{h,j}=[z^h]\prod_{i=1}^r(1-B_i z)^{j/r},
             \qquad 1\le j\le h-1.
\end{equation}
Then
\begin{equation}\label{sd:eq:localderivative}
 \boxed{\displaystyle
 \mathcal L_W^{(m)}(0)=
 \sum_{h\ge2}p_h(-1)^{m+1}m!r^m h!
  \sum_{j=1}^{h-1}
    \frac{(-1)^jC_{h,j}}{j!(h-1-j)!\,j^{m+1}}.}
\end{equation}
The coefficient extraction in \eqref{sd:eq:Ccoeff} is finite and
algebraic; it makes no branch choice beyond the formal series with
constant coefficient one.
\end{theorem}
\begin{proof}
Dirichlet moments give
\[
 \mathbb E_{r,s}\left(\sum_i t_iB_i\right)^h
  =\frac{h!}{(rs)_h}
       \sum_{h_1+\cdots+h_r=h}
                    \prod_i\frac{(s)_{h_i}B_i^{h_i}}{h_i!}.
\]
This proves \eqref{sd:eq:local}.  At $s=0$ its continuation is the
average of the vertex values: only compositions supported on a single
index survive after the common factor $s$ has been cancelled.

For the derivative formula put $v=rs$.  The $h$th moment becomes
\[
 M_h(v)=\frac{h!}{(v)_h}
                   [z^h]\prod_i(1-B_i z)^{-v/r}.
\]
After cancellation at zero, it is a rational function of numerator and
denominator degree at most $h-1$.  Its polynomial part is constant.
Its possible poles are simple and lie at $v=-j$, $1\le j<h$, with
residues
\[
 \operatorname{Res}_{v=-j}M_h(v)
   =\frac{(-1)^j h!C_{h,j}}{j!(h-1-j)!}.
\]
Differentiate the partial fractions $m$ times at $v=rs=0$ and subtract
from the vertex average.  This is exactly
\eqref{sd:eq:localderivative}.
\end{proof}

For example, let $V(\boldsymbol B)=\sum_{i<j}(B_i-B_j)^2$.  For the
first nontrivial odd weight,
\begin{equation}\label{sd:eq:cubiclocal}
 \mathcal L_{x^3}(s)=\frac{sV(\boldsymbol B)}{4r(1+rs)},\qquad
 \mathcal L_{x^3}^{(m)}(0)
       =\frac{m!(-r)^{m-1}}{4r}V(\boldsymbol B).
\end{equation}
More generally, after division by $m!$, the sequence in
\eqref{sd:eq:localderivative} is a finite linear combination of powers
$(-r/j)^m$.  It therefore satisfies an explicit finite recurrence in
$m$.  This rational contribution is the part supplied by the zero
one-mass jet.  It should not be read as saying that Stieltjes or residue
terms inside the lower jets $F_1,F_2,\ldots$ disappear.

\subsection{Every derivative uses only finitely many lower jets}
\label{sd:subsec:allorders}

For $k\ge2$ and $q\ge0$ define
\begin{equation}\label{sd:eq:P}
 P_{k,q}(L)=[z^q]\exp\left(
 Lz+\sum_{n\ge2}\frac{(-1)^n\zeta(n)}n(k^n-k)z^n\right).
\end{equation}
Only finitely many terms of the exponential are required at any fixed
$q$.  In particular,
\[
 P_{k,0}=1,\qquad P_{k,1}=L,\qquad
 P_{k,2}=\frac{L^2}{2}+\frac{k^2-k}{2}\zeta(2).
\]
Put $\ell_S(t)=\sum_{i\in S}\log t_i$.

\begin{theorem}[All-order formula and face-size restriction]
\label{sd:thm:allorders}
For every integer $m\ge1$,
\begin{equation}\label{sd:eq:allorders}
 \boxed{\begin{aligned}
 \Delta_m(W)={}&\mathcal L_W^{(m)}(0)\\
 &-\frac{m!}{r}
  \sum_{k=2}^{\min(r,m)}\ \sum_{|S|=k}
  \sum_{j=1}^{m-k+1}\frac{r^j}{j!}
   \int_{\Sigma_S}\mathcal H_S(F_j;t)
       P_{k,m-k+1-j}(\ell_S(t))
                     \frac{dt}{\prod_{i\in S}t_i}.
 \end{aligned}}
\end{equation}
Every displayed simplex integral is absolutely convergent.  After the
explicit zero-jet term is removed, a $k$-vertex face starts at order
$s^k$ and uses one-mass derivatives of order at most $m-k+1$ in the
$m$th derivative.
\end{theorem}
\begin{proof}
The classical Taylor series of $\log\Gamma(1+z)$ gives
\[
 Q_k(s)=\exp\left(
       \sum_{n\ge2}\frac{(-1)^n\zeta(n)}n(k^n-k)s^n\right);
\]
the Euler-constant terms cancel.  Once $F_0$ has been removed, the
inner function in \eqref{sd:eq:master} is
$\sum_{j\ge1}r^js^jF_j/j!$.  Multiplication by $s^{k-1}$ proves
that this face starts at degree $k$.  Extract the coefficient of $s^m$
from \eqref{sd:eq:master}; the remaining exponential is precisely
\eqref{sd:eq:P}.  Lemma~\ref{sd:lem:boundary} justifies extraction under
each integral.  Multiplication by $m!$ gives the theorem.
\end{proof}

For a compact statement of the first instances write
\[
 I_S(F)=\int_{\Sigma_S}\mathcal H_S(F;t)
                       \frac{dt}{\prod_i t_i},\qquad
 J_S(F)=\int_{\Sigma_S}\mathcal H_S(F;t)\ell_S(t)
                       \frac{dt}{\prod_i t_i},
 \qquad L_m=\mathcal L_W^{(m)}(0).
\]
All subset sums below use distinct labels.  The complete second and
third derivatives are
\begin{align}
 \Delta_2(W)&=L_2-2\sum_{|S|=2}I_S(F_1),
                    \label{sd:eq:second}\\
 \Delta_3(W)&=L_3-3r\sum_{|S|=2}I_S(F_2)
                  -6\sum_{|S|=2}J_S(F_1)
                  -6\sum_{|S|=3}I_S(F_1).
                    \label{sd:eq:third}
\end{align}
The fourth derivative is
\begin{align}\label{sd:eq:fourth}
 \Delta_4(W)={}&L_4-4r^2\sum_{|S|=2}I_S(F_3)
                    -12r\sum_{|S|=2}J_S(F_2)\notag\\
 &-24\sum_{|S|=2}\int_{\Sigma_S}\mathcal H_S(F_1;t)
             \left(\frac{\ell_S(t)^2}{2}+\zeta(2)\right)
                                  \frac{dt}{\prod_i t_i}\notag\\
 &-12r\sum_{|S|=3}I_S(F_2)-24\sum_{|S|=3}J_S(F_1)
                              -24\sum_{|S|=4}I_S(F_1).
\end{align}
A subset sum is empty when its size exceeds $r$.  The term $\zeta(2)$
in the pair integral of \eqref{sd:eq:fourth} comes from $Q_2$ and is
part of the normalization.  In particular the third derivative needs
only pair integrals of the second one-mass jet and pair and triple
integrals of the first jet.

\subsection{Explicit Gamma integrands and a nonzero face family}
\label{sd:subsec:gamma}

For $W=1$, first-derivative factorization gives the classical expression
\begin{equation}\label{sd:eq:gamma}
 F_1(B)=\log\Gamma(a+i\sqrt B)+\log\Gamma(a-i\sqrt B)
                                      -\log(2\pi),
\end{equation}
using the analytic logarithm of Gamma in the right half-plane.
The corresponding even-weight expression, inherited from
first-derivative factorization, is
\begin{equation}\label{sd:eq:evenhurwitz}
 F_1(B)=\sum_{\varepsilon=\pm1}\sum_{j=0}^{2d}
   \binom{2d}{j}(-\varepsilon i\sqrt B)^{2d-j}
               \zeta'(-j,a+\varepsilon i\sqrt B),
             \qquad W=x^{2d}.
\end{equation}
Thus the largest face size $k=m$ in \eqref{sd:eq:allorders} requires
only ordinary Gamma at $W=1$, and a finite Hurwitz-derivative expression
for every even polynomial weight.  Its coefficients and integration
kernels are fully explicit.

The following identity shows that the higher faces are substantive
terms of this particular representation.

\begin{proposition}[A factorial identity for every face size]
\label{sd:prop:quadraticface}
For $k=|S|\ge2$,
\begin{equation}\label{sd:eq:quadraticface}
 \boxed{\displaystyle
 I_S(B\mapsto B^2)=(-1)^{k-1}(k-2)!
                             \sum_{\{i,j\}\subseteq S}(B_i-B_j)^2.}
\end{equation}
Consequently, if $B_i=B+\epsilon\beta_i$ and $F$ is $C^3$ near $B$,
\begin{equation}\label{sd:eq:collisionleading}
 I_S(F)=\frac{\epsilon^2F''(B)}2(-1)^{k-1}(k-2)!
           \sum_{\{i,j\}\subseteq S}(\beta_i-\beta_j)^2
                         +O(\epsilon^3).
\end{equation}
\end{proposition}
\begin{proof}
On the nonnegative coordinate cone the quadratic perspective has the
variance decomposition
\[
 t_T B_T^2=\sum_{i\in T}t_iB_i^2
       -\sum_{\{i,j\}\subseteq T}
                 \frac{t_it_j}{t_T}(B_i-B_j)^2.
\]
The first sum is annihilated by the face subtraction.  Writing
$t_T^{-1}=\int_0^\infty e^{-vt_T}\,dv$ in the second gives
\begin{align*}
 \mathcal H_S(B^2;t)
  ={}&(-1)^{k-1}\sum_{\{i,j\}\subseteq S}(B_i-B_j)^2t_it_j\\
 &\qquad\times\int_0^\infty e^{-v(t_i+t_j)}
                 \prod_{\ell\in S\setminus\{i,j\}}
                         (1-e^{-vt_\ell})\,dv.
\end{align*}
Since $\mathcal H_S$ is homogeneous of degree one, radial integration
shows that $I_S$ equals the integral of
$e^{-\sum_i t_i}\mathcal H_S(t)/\prod_i t_i$ on the positive
orthant.  The radial factor is $\int_0^\infty e^{-\rho}\,d\rho=1$.
After the displayed sign is removed, every term is nonnegative, so
Tonelli's theorem applies.  The two distinguished coordinates give
$(1+v)^{-2}$, and each other coordinate gives
\[
 \int_0^\infty e^{-t}\frac{1-e^{-vt}}{t}\,dt=\log(1+v).
\]
The remaining integral is
\[
 \int_0^\infty\frac{\log^{k-2}(1+v)}{(1+v)^2}\,dv=(k-2)!.
\]
This proves \eqref{sd:eq:quadraticface}.  Taylor-expand $F(B+\epsilon
\beta)$ through degree two.  The constant and affine parts are
annihilated; Lemma~\ref{sd:lem:boundary} controls the integrated
remainder uniformly, proving \eqref{sd:eq:collisionleading}.
\end{proof}

For \eqref{sd:eq:gamma},
\[
 F_1''(B)=\sum_{n\ge0}\frac1{(x_n^2+B)^2}>0.
\]
Hence every face size has a nonzero quadratic collision coefficient
when the perturbations are not all equal.  This does not prove
irreducibility into any other chosen function class.  In particular,
$\mathcal H_S$ need not vanish merely because two of several masses
coincide; ``face contribution'' is the precise description used here.

\subsection{Periodic weights and cancellation of the mean}
\label{sd:subsec:periodic}

\begin{corollary}[Mean-zero periodic weights remove the zero-jet term]
\label{sd:cor:periodic}
Let $c_n$ be $q$-periodic, and insert $c_n$ in every spectral sum.
Theorems~\ref{sd:thm:master} and \ref{sd:thm:allorders} still hold.
The rational zero-jet contribution is
\begin{equation}\label{sd:eq:periodiclocal}
 \mathcal L_{c,W}(s)=\overline c\,\mathcal L_W(s),\qquad
                  \overline c=\frac1q\sum_{v=0}^{q-1}c_v.
\end{equation}
Thus $\overline c=0$ removes this term for \emph{every} polynomial
weight, including odd weights.
\end{corollary}
\begin{proof}
The auxiliary Dirichlet series is
\[
 \sum_{n\ge0}c_n(n+a)^{-z}
   =q^{-z}\sum_{v=0}^{q-1}c_v
                           \zeta\!\left(z,\frac{a+v}{q}\right).
\]
Its residue at $z=1$ is $\overline c$.  Repeat the proof of
Lemma~\ref{sd:lem:zero}; the zero jet becomes
\[
 F_{c,0}(B)=\sum_d w_dq^d
       \sum_v c_v\zeta\!\left(-d,\frac{a+v}{q}\right)
                    +\overline c\sum_{h\ge1}p_hB^h.
\]
The mass-independent part is annihilated by every face subtraction.
All other steps use linearity and normal convergence; periodic weights
are bounded and cause no additional convergence issue.
\end{proof}

For $W=1$ the first jet has the explicit finite dictionary
\begin{align}\label{sd:eq:periodicgamma}
 F_{c,1}(B)=\sum_{v=0}^{q-1}c_v\bigg\{&
 \log\Gamma\!\left(\frac{a+v+i\sqrt B}{q}\right)
 +\log\Gamma\!\left(\frac{a+v-i\sqrt B}{q}\right)\notag\\
 &-\log(2\pi)-2\log q
                         \left(\frac12-\frac{a+v}{q}\right)\bigg\}.
\end{align}
This follows by scaling each residue class to
$q^{-2s}E_1(s;B/q^2)$ with shift $(a+v)/q$, where the subscript
$1$ here denotes the polynomial weight.  Formula
\eqref{sd:eq:periodicgamma} therefore supplies ordinary Gamma
integrands for alternating weights and character projections.  The
mass-independent constants can be omitted inside $\mathcal H_S$,
but are retained in the full first-jet identity.

\subsection{Independent evaluation and research consequences}
\label{sd:subsec:validation}

For $a=2$, $W=1$, $r=3$, and masses
$(B_1,B_2,B_3)=(0.2,0.7,1.3)$, the directly evaluated spectral derivative is
\begin{equation}\label{sd:eq:numeric}
 \Delta_3=-1.8999673234892670916122436912924709222747240234\ldots.
\end{equation}
One calculation used a finite ordinary spectral sum and an independently
summed binomial--Hurwitz tail.  Changing the head length and tail degree
from $(55,13)$ to $(85,17)$ changed the result by $3.04\times10^{-46}$.
A second calculation evaluated \eqref{sd:eq:third}: the pair terms used
convergent one-mass expansions, while the triple integral used Gamma
values from \eqref{sd:eq:gamma}.  The errors in the full face formula
at Gauss--Legendre orders $16$, $24$, and $32$, with $45$-digit
arithmetic, were respectively
$1.62\times10^{-18}$, $2.54\times10^{-27}$, and
$8.20\times10^{-32}$.  The last error is limited by the retained
one-mass expansion.  The triple integral was
\[
 I_{\{1,2,3\}}(F_1)
    =0.0564600947235384049461537670929012515\ldots.
\]
For stable numerical integration the simplex was divided into its six
coordinate orderings.  On each, use
$t=(1,u,uv)/(1+u+uv)$; the weighted measure becomes
$dt/(t_1t_2t_3)=du\,dv/(uv)$, and the face cancellation removes the
corner singularity.  The code and precision-dependent outputs are
included with the article.  These are numerical diagnostics, not
interval enclosures or replacements for the proofs.

The results close the explicit face-subtraction problem and supply a
complete odd-weight zero-jet correction in this lower-jet framework.
Further concrete questions are to evaluate special Gamma face integrals
in a fixed polylogarithmic function class; to obtain comparable
multivariable formulas when the spectral orders of the factors vary
independently; and to derive closed all-order collision coefficients for
the individual face integrals beyond \eqref{sd:eq:collisionleading}.
The theorem itself does not require such reductions and makes no claim
that the unresolved Gaussian period conjectures are settled.
